\documentclass[11pt]{article}

\usepackage[margin=1in]{geometry}
\usepackage{amsmath,amssymb,amsfonts}
\usepackage{booktabs}
\usepackage{graphicx}
\usepackage{xcolor}
\usepackage{microtype}
\usepackage{placeins}
\usepackage{caption}
\usepackage{enumitem}
\usepackage{hyperref}
\usepackage[numbers,sort&compress]{natbib}

% Keep all body paragraphs flush with the left margin, including in generated builds.
\AtBeginDocument{%
  \setlength{\parindent}{0pt}%
  \setlength{\parskip}{0.48em}%
}

\hypersetup{
  colorlinks=true,
  linkcolor=blue!55!black,
  citecolor=blue!55!black,
  urlcolor=blue!55!black
}

% Full-width abstract with controlled paragraph alignment. The standard article
% abstract uses a narrower quotation block, which looked visually off-center in
% the compact arXiv build.
\renewenvironment{abstract}{%
  \par\small
  \begin{center}\bfseries Abstract\end{center}%
  \vspace{-0.45em}%
  \setlength{\parindent}{0pt}%
  \setlength{\parskip}{0.35em}%
  \ignorespaces
}{%
  \par\normalsize
}

\title{When Site Signal and Biology Collide:\\
PathoAlign Neural Identifiability and Dominant-Site Stress Tests for Computational Pathology}

\author{Matthew Vaishnav\\
Independent Researcher}

\date{\today}

\begin{document}
\maketitle

\begin{abstract}
Computational-pathology representations often contain both clinically useful biological information and acquisition-dependent information associated with staining, scanners, preparation protocols, institutions, and annotation workflows. These components can be statistically entangled. A model that removes site information too aggressively may therefore erase diagnostic morphology, while a federated model that assigns influence only by sample count may amplify a site-specific training process that is poorly aligned with the declared validation objective.

I study these two related failure modes through PathoAlign, a neural identifiability program for separating biological and acquisition factors, and through dominant-site federated stress tests over PANDA-derived Phikon features. PathoAlign experiments include task-only neural networks, gradient-reversal site-adversarial models, factorized diagnostic/site branches, pair-consistency objectives, learned acquisition operators, and hybrid curricula. A locked confirmatory experiment found no task-accuracy advantage for the adversarial PathoAlign model over the tumor-only baseline, but did find a small reproducible reduction in site leakage. Factorized modeling selected the diagnostic-only baseline, and unconditional site-component subtraction was falsified because the learned site component remained strongly tumor-predictive. These negative results establish that center invariance alone is not sufficient and that nuisance removal must preserve biology conditionally.

Paired biology-preserving acquisition interventions provide a stronger information source: differencing paired representations cancels the shared biological term and exposes acquisition variation. Controlled synthetic experiments vary observational sample size, paired-anchor count, biological--acquisition overlap, linear versus nonlinear mixing, method, and seed. The resulting phase maps reveal a two-resource recovery boundary. At the primary threshold, an interval/right-censored fit gives a local empirical relation in which more observational data reduces the required paired-anchor count, while overlap and nonlinear mixing increase it. This relation is treated as a falsifiable synthetic scaling hypothesis, not a universal law.

The federated experiments study a complementary influence problem. A fixed detector-switch rule calibrated under dominant-site label-noise stress transfers to conservative ordinal threshold shift, while maintaining low clean-regime switching and producing positive improvements at stronger shift levels. Together, the results support a narrow conclusion: site variation cannot be treated as pure nuisance, sample volume cannot be treated as complete aggregation authority, and pathology models require auditable tests of what information is removed, retained, and allowed to dominate.
\end{abstract}

\section{Introduction}
Federated learning is attractive in medicine because raw patient data is difficult, risky, and often impossible to centralize. In pathology, the motivation is especially clear: whole-slide images and derived feature representations may be governed by institutional policy, patient privacy requirements, storage constraints, and clinical governance. Federated learning offers a way for multiple sites to train collaboratively while keeping raw local data inside each environment.

However, federated learning does not remove the question of influence. It changes the question from ``who owns the data?'' to ``how much should each site shape the shared model?'' Standard FedAvg answers that question with a simple rule: clients with more samples receive more aggregation influence \citep{mcmahan2017communication}.

A related problem appears inside the representation itself. Pathology features may encode tumor morphology together with staining, scanner, preparation, annotation, and institutional effects. If those components overlap, an unconditional site-removal operation can remove clinically useful information along with acquisition variation. The central PathoAlign question is therefore not merely whether a representation predicts site, but whether the model can identify which site-associated directions are nuisance and which are carrying biological signal.

I study controlled versions of both problems. In the federated setting, I simulate multi-site learning over PANDA-derived Phikon feature vectors and perturb the largest simulated site's training labels while keeping validation labels clean. In the PathoAlign setting, I train neural models under known synthetic biological and acquisition factors, test adversarial and factorized representations, introduce paired acquisition interventions that preserve biology, and measure where recovery succeeds or fails.

\subsection*{Working thesis}
In computational pathology, site information is neither automatically useless nor automatically trustworthy. It can contain nuisance variation, biological signal, or both. Safe modeling therefore requires conditional, tumor-preserving nuisance removal at the representation level and auditable limits on sample-size dominance at the federation level.

% The introductory architecture figure was intentionally retired; do not reinsert federated_pathology_pipeline_diagram.tex.
\paragraph{Method at a glance.}
For client updates $\Delta\theta_k$, the two candidate global updates are
\[
\Delta\theta_{\mathrm{FA}}
=\sum_{k=1}^{K}\frac{n_k}{\sum_{j=1}^{K}n_j}\Delta\theta_k,
\qquad
\Delta\theta_{\mathrm{DA}}
=\sum_{k=1}^{K}\widetilde w_k\Delta\theta_k.
\]
A clean-calibrated diagnostic vector
\[
d(r)=\bigl(\kappa_{\mathrm{global}},\kappa_{\min},S_\kappa,
\operatorname{MAOE},R_{\mathrm{sev}}\bigr)
\]
induces the trigger
\[
C(r)=\sum_{j=1}^{5}I_j(r),
\qquad
T(r)=\mathbf{1}[C(r)\geq3],
\]
and the selected policy is
\[
M_{\pi}(r)
=[1-T(r)]M_{\mathrm{FA}}(r)+T(r)M_{\mathrm{DA}}(r).
\]
At the representation level, paired acquisitions of the same biological specimen satisfy
\[
x_i=Bz_{b,i}+Az_{a,i}+\varepsilon_i,
\qquad
x_i'=Bz_{b,i}+Az'_{a,i}+\varepsilon_i',
\]
so that
\[
\Delta x_i=x_i'-x_i
=A(z'_{a,i}-z_{a,i})+(\varepsilon_i'-\varepsilon_i).
\]
The shared biological term cancels, providing direct supervision for acquisition-sensitive directions while biological representations are trained to remain stable.

\section{Contributions}
I make eight main contributions.

\begin{enumerate}
  \item \textbf{A neural identifiability framework for computational pathology.} PathoAlign treats biological and acquisition signals as potentially entangled latent factors and evaluates task-only, gradient-reversal adversarial, factorized, pair-consistency, learned-operator, and hybrid neural models under controlled conditions.

  \item \textbf{A locked neural-network confirmation and falsification result.} The adversarial PathoAlign model reduced site leakage slightly but did not improve held-out tumor accuracy over the tumor-only baseline. A factorized model selected the diagnostic-only baseline, and unconditional site-component subtraction failed because the removed component remained tumor-predictive. These results rule out the simple claim that stronger center invariance is automatically better.

  \item \textbf{A paired-intervention identifiability mechanism.} Biology-preserving acquisition pairs cancel the shared biological term under differencing, turning otherwise ambiguous site variation into a learnable intervention signal and motivating pair-consistency and operator-based neural objectives.

  \item \textbf{Failure and recovery phase maps.} I vary observational sample size, paired-anchor count, biological--acquisition overlap, nonlinear mixing, neural method, and seed to identify regimes where factor separation is possible, unstable, or right-censored beyond the tested intervention budget.

  \item \textbf{A two-resource empirical recovery law.} Interval/right-censored analysis links the required paired-anchor count to observational sample size, overlap, and nonlinear mixing. The relationship is reported as a threshold-dependent local empirical family and is reserved for held-out falsification rather than presented as a universal theorem.

  \item \textbf{A sample-volume / site-signal alignment failure-mode framing.} I reframe dominant-client failure not as institutional reliability ranking, but as an audit of a modeling assumption that FedAvg already makes silently: larger client equals more influence.

  \item \textbf{A fixed detector-switch transfer result in neural federated pathology.} Compact slide-level neural networks trained over PANDA-derived Phikon representations show that a detector calibrated under dominant-site label-noise stress can transfer to conservative ordinal threshold shift with low clean-regime switching and positive gains at stronger shift levels.

  \item \textbf{Robustness, negative-result, and ethical reporting.} The project includes locked confirmation, paired bootstrap analysis, threshold sensitivity, leave-one-diagnostic-family-out ablation, calibration sensitivity, explicit unsupported-claim boundaries, and a non-blameworthy interpretation of site variation.
\end{enumerate}

\section{PathoAlign neural identifiability program}
\subsection{Neural models}
PathoAlign is not a single network architecture. It is a controlled neural experimental program for testing what forms of supervision can separate biological and acquisition factors without discarding diagnostic information. The initial model used a shared feature encoder with a tumor-prediction head and a gradient-reversal site head. The adversarial objective encouraged tumor prediction while discouraging directly recoverable center information.

A later factorized model used a shared stem with separate diagnostic and site branches, reconstruction pressure, and a covariance penalty intended to reduce redundancy between the branches. Subsequent synthetic experiments added pair-consistency models that force biology-preserving pairs to agree in biological representation, learned operators that model acquisition-induced transformations, and hybrid curricula that combine observational task supervision with paired-intervention supervision.

The methods are compared using more than task accuracy. Evaluation includes site leakage from biological representations, tumor predictability from acquisition representations, factor separation, paired biological invariance, counterfactual transport, canonical-correlation recovery, and Procrustes alignment to known latent factors in the synthetic setting.

\subsection{Locked confirmation and negative results}
Development experiments initially suggested a small tumor-accuracy improvement and reduced site leakage for the adversarial model. A locked confirmatory run over new seeds did not reproduce the accuracy advantage: PathoAlign achieved approximately 0.9264 test accuracy versus 0.9270 for the tumor-only baseline, with a paired difference of approximately $-0.000625$ and an interval crossing zero. The site-probe difference was approximately $-0.003258$, with the paired bootstrap interval excluding zero, supporting a small reduction in recoverable site information.

This result is scientifically useful precisely because it is not a task-performance win. It shows that the adversarial neural objective can alter representation leakage without establishing better diagnosis. The factorized v1 model likewise selected the diagnostic-only baseline under validation-based model selection. A later unconditional subtraction test was more decisive: the learned site component was itself strongly tumor-predictive, so subtracting it removed useful biology. The resulting conclusion is that center invariance alone is not sufficient and that site-associated information must be decomposed conditionally rather than erased wholesale.

\subsection{Why paired interventions change the problem}
Observational data alone does not identify whether a site-associated direction is nuisance, biology, or an entangled mixture. Biology-preserving acquisition pairs add intervention structure: the same specimen or latent biology is observed under different acquisition states. Under the controlled generative model, subtraction cancels the shared biology and exposes acquisition variation. This changes the problem from unsupervised nuisance guessing to partially supervised identifiability.

The synthetic PathoAlign experiments therefore treat unpaired observations and paired interventions as distinct resources. Observational data teaches the task manifold, while paired interventions reveal which changes should not alter biological representation. The resulting phase maps test how much of each resource is required as overlap and nonlinearity increase.

\section{Ethical framing: site-signal alignment, not institutional worth}
This work does \emph{not} claim that some hospitals, pathologists, or institutions are inherently more reliable, competent, or trustworthy than others.

The term \emph{site-signal alignment} is used in a narrow modeling sense: whether a simulated client's training signal appears aligned with the declared validation objective under a specific experimental setup. A client may appear misaligned for many non-blameworthy reasons, including different grading thresholds, staining or scanning protocol differences, local case mix, patient-population differences, annotation workflow differences, label-source differences, historical reporting practice, local clinical policy, or pathologist disagreement.

Dominance-aware aggregation should therefore not be interpreted as an institutional ranking system. It is an audit mechanism for a modeling assumption that FedAvg already makes silently: larger client equals more influence.

The ethical purpose of this work is to make that assumption visible, stress-testable, and contestable. Any real deployment would require governance, local clinical review, pathologist input, bias auditing, institutional agreement on validation objectives, prospective validation, security review, and regulatory review.

\section{Experimental setup}
\subsection{Dataset and representation}
The federated experiments use PANDA-derived pathology feature representations. Whole-slide images are represented through precomputed Phikon features. Patch-level representations are pooled into one standardized slide vector $x\in\mathbb{R}^{D}$ for the downstream six-class ISUP grading task. The PathoAlign identifiability experiments use both controlled synthetic representations with known biological and acquisition factors and pathology-derived representations for leakage and task-retention tests.

\subsection{Slide-level neural network}
Every federated client trains the same neural network architecture, so comparisons isolate aggregation policy rather than model-capacity differences. The slide classifier is a compact multilayer perceptron
\begin{equation}
  h_1=\operatorname{Dropout}_{0.20}\!\left(\operatorname{ReLU}(W_1x+b_1)\right),
\end{equation}
\begin{equation}
  h_2=\operatorname{Dropout}_{0.20}\!\left(\operatorname{ReLU}(W_2h_1+b_2)\right),
\end{equation}
\begin{equation}
  \widehat y=W_3h_2+b_3,
\end{equation}
where $W_1$ maps the pooled Phikon input from $D$ dimensions to 256 hidden units, $W_2$ maps 256 units to 128, and $W_3$ produces six logits corresponding to ISUP grades 0--5. Softmax probabilities are used for prediction and evaluation, while local optimization uses multiclass cross-entropy.

Unless otherwise specified by an experiment configuration, each simulated client performs one local epoch per federated round using AdamW with learning rate $10^{-3}$, weight decay $10^{-4}$, and batch size 64. The global experiment runs for five communication rounds. Each strategy starts from the same seeded architecture and uses the same local optimizer; only the server-side aggregation weights or detector-controlled switching policy differ.

This network is deliberately modest. The scientific target is the relationship between client influence and site-signal alignment, not a claim that this MLP is a new pathology backbone. Phikon supplies the pretrained pathology representation, and the federated neural network learns the task-specific slide classifier over those representations.

\subsection{Simulated federations}
The feature dataset is partitioned into simulated clients. One client is designated as the dominant client by sample volume. The clean validation objective is held fixed while the dominant client's training labels are perturbed.

\subsection{Stress modes}
Two dominant-site stress modes are considered:
\begin{enumerate}
  \item random label corruption at the dominant site;
  \item systematic ordinal threshold shift at the dominant site.
\end{enumerate}

\subsection{Evaluation metrics}
The primary federated metrics are global QWK, macro-F1, and worst-site QWK. Detector diagnostics additionally include site-specific QWK spread, mean absolute ordinal error, severe-error rate, and prediction entropy where available. PathoAlign metrics additionally audit tumor retention, site leakage, cross-factor leakage, paired invariance, transport error, and latent-factor recovery.

\section{Mathematical notation}
Let there be $K$ simulated clients. Client $k$ has $n_k$ local training samples and produces a local model update $\Delta\theta_k$. Standard FedAvg assigns aggregation weight
\begin{equation}
  w_k=\frac{n_k}{\sum_{j=1}^{K}n_j}
\end{equation}
and forms the global update
\begin{equation}
  \Delta\theta=\sum_{k=1}^{K}w_k\Delta\theta_k.
\end{equation}

The concern studied here is that $n_k$ is a volume term, not an alignment term. A high-volume client may have a training-label process that is less aligned with the declared validation objective.

Let $A_k$ denote the task-specific alignment of client $k$'s training signal with the declared validation objective. The claim is not that $A_k$ is an institutional property. It is a run-specific modeling quantity. The failure mode appears when the dominant sample-volume client has high $n_k$ but reduced $A_k$.

Let $d_1,\ldots,d_m$ be clean-calibrated validation diagnostics. Each diagnostic has a clean safe range defined by lower and/or upper quantiles. For a run $r$, define
\begin{equation}
  I_i(r)=
  \begin{cases}
    1, & \text{if diagnostic }d_i\text{ leaves its clean-calibrated safe range},\\
    0, & \text{otherwise.}
  \end{cases}
\end{equation}

The fixed detector count is
\begin{equation}
  C(r)=\sum_{i=1}^{m}I_i(r).
\end{equation}

The fixed transferred rule uses five diagnostics and switches when at least three are abnormal:
\begin{equation}
  T(r)=\mathbf{1}[C(r)\geq3].
\end{equation}

Let $M_{\mathrm{FA}}(r)$ be the metric under FedAvg and $M_{\mathrm{DA}}(r)$ the metric under the selected dominance-aware alternative. The detector-controlled metric is
\begin{equation}
  M_{\pi}(r)=[1-T(r)]M_{\mathrm{FA}}(r)+T(r)M_{\mathrm{DA}}(r).
\end{equation}

Prediction entropy is available but is not used in the fixed headline rule.

\section{Fixed detector transfer result}
The detector was calibrated under dominant-site label-noise stress and then transferred without retuning to conservative ordinal threshold shift.

At 35\% and 45\% conservative shift, the fixed rule produced positive paired improvements in global QWK, macro-F1, and worst-site QWK while keeping clean-regime switching at 13.3\%. The result is therefore not a claim that the alternative aggregation rule is universally superior. It is a claim that conditional switching can reduce risk when a detector identifies evidence that sample-size dominance has become unsafe.

\section{Diagnostic interpretation}
The detector is not driven by one heuristic alone. Across the transferred threshold-shift runs, the most frequent trigger component is mean absolute ordinal error increase, followed by degradation in global and worst-site QWK. Site-QWK spread contributes, but it is not the sole or dominant explanation for the result.

This matters because a detector based only on cross-site dispersion could confuse ordinary site heterogeneity with harmful dominant-site misalignment. The transferred rule instead combines global degradation, worst-site degradation, dispersion, and ordinal-error severity.

\section{Ablation and calibration sensitivity}
Leave-one-diagnostic-family-out ablation shows that removing the most frequent diagnostic family, mean absolute ordinal error, reduces trigger rate but does not collapse the positive 35\% / 45\% transfer result.

Calibration-sensitivity analysis varies the lower quantile, upper quantile, and minimum trigger count around the fixed rule. Of 36 nearby configurations, 29 preserve clean switching at or below 20\% and positive global-QWK, macro-F1, and worst-site-QWK deltas at both 35\% and 45\% conservative shift.

\section{Limitations}
The experiments include controlled synthetic generators and simulated federated clients over fixed feature representations rather than prospective multi-hospital deployment. The shared federated task model is a compact slide-level MLP on pooled Phikon vectors; results may change with attention-based MIL, transformer-based MIL, foundation-model fine-tuning, or full patch-level training. The PathoAlign scaling relation is local to the tested synthetic regime, threshold definition, pair-count grid, and methods. The locked adversarial result reduces leakage but does not improve task accuracy. The stress modes are controlled perturbations, not prospective clinical observations. The detector is calibrated from clean runs and depends on the chosen validation objective. The work does not establish clinical safety, regulatory readiness, fairness across patient subgroups, resistance to adversarial clients, privacy guarantees beyond local data retention, or performance under real deployment constraints.

\section{Discussion}
The combined results support a broader but still bounded conclusion: site-associated information must be audited rather than automatically removed or automatically trusted. At the representation level, site information can be entangled with diagnostic morphology, making unconditional subtraction unsafe. At the federation level, sample size can overstate the authority of a training signal that is poorly aligned with the declared objective.

Paired interventions provide a principled way to distinguish changes that should alter acquisition representation from features that should remain biologically stable. The synthetic two-resource phase map suggests that ordinary observations and paired interventions play different roles: observations estimate the task manifold, while interventions reveal transformation structure. The required intervention budget rises sharply with biological--acquisition overlap and nonlinear mixing.

The appropriate response is not to distrust large sites or to force complete site invariance. It is to make both representation removal and aggregation influence auditable. Future work should test the frozen PathoAlign recovery predictions on unseen conditions, evaluate conditional tumor-preserving low-rank removal on real pathology data, extend the models to MIL and foundation-model feature bags, and validate the federation mechanism under real institutional partitions and governance.

\section{Conclusion}
PathoAlign shows that neural site invariance is not equivalent to biological identifiability. A model can reduce site leakage without improving diagnosis, and a learned site component can still carry tumor information. Paired biology-preserving interventions change the information structure and support a measurable recovery boundary linking observational data, intervention count, overlap, and nonlinearity. Complementary federated experiments show that sample-size dominance should also be treated as an auditable assumption. Together, these results motivate pathology systems that explicitly test what information is removed, retained, and allowed to dominate rather than assuming that site signal is pure nuisance or that more data automatically deserves more influence.

\bibliographystyle{plainnat}
\bibliography{references}

\end{document}
